% *{{ preamble
\documentclass[12pt, pdf, hyperref={draft}, usenames, dvipsnames]{beamer}
\usepackage{libertine, textcomp, amssymb, wasysym, color, ulem, verbatim}
\usepackage{floatrow, subcaption}
\usepackage[makeroom]{cancel}
\usepackage{mathtools}
% hide controls
\usenavigationsymbolstemplate{}
\usepackage[T1]{fontenc}
\synctex=1
\setbeamertemplate{caption}[numbered]
% }}*

% *{{ item colouring and themes
\usecolortheme{beaver}
\setbeamertemplate{footline}[frame number]
\definecolor{LancsRed}{HTML}{B5121B}
\setbeamercolor{item}{fg=LancsRed}
\setbeamercolor{title}{fg=LancsRed}
\setbeamercolor{structure}{fg=LancsRed}
\setbeamercolor{frametitle}{fg=LancsRed}
\setbeamercolor{footnote mark}{fg=LancsRed}
\setbeamercolor{footnote}{fg=LancsRed}
\setbeamercolor{bibliography entry author}{fg=LancsRed}
\setbeamercolor{bibliography entry title}{fg=black}
\setbeamercolor{bibliography entry note}{fg=black}
\setbeamercolor{section in toc}{fg=LancsRed}
%\setbeamercolor{subsection in toc}{fg=gray}
\newcommand{\gitem}[1]{\setbeamercolor{item}{fg=ForestGreen}\item[$\checkmark$] #1}
\newcommand{\bitem}[1]{\setbeamercolor{item}{fg=LancsRed}\item[$\times$] #1}
\newcommand{\nitem}[1]{\setbeamercolor{item}{fg=NavyBlue}\item[$-$] #1}
% }}*

% *{{ niceties
\newcommand{\dd}{\mathrm{d}}
\newcommand{\e}{\mathrm{e}}
\newcommand{\ket}[1]{\lvert{#1}\rangle}
\newcommand{\bra}[1]{\langle{#1}\rvert}
\newcommand{\expt}[1]{\langle{#1}\rangle}
\newcommand{\red}[1]{{\bf\color{LancsRed}{#1}}}
\newcommand{\blue}[1]{{\bf\color{NavyBlue}{#1}}}
\newcommand{\green}[1]{{\bf\color{ForestGreen}{#1}}}

% }}*

% *{{ title, subtitle, inst., date
\title{Modelling Charge Complexes in \\
       2D Materials \& Their Heterostructures}
\subtitle{Northern Quantum Meeting II}

% authors and institutions
%\author{\emph{Ryan J. Hunt}\inst{1},
%        Neil D. Drummond\inst{1} \\
%        and Vladimir I. Fal'ko\inst{2}}

\author{Ryan Hunt \\
15$^{\text{th}}$ December, 2017}

%\institute[]{
%
%  \inst{1}
%  Department of Physics,\\
%  Lancaster University
%  \and
%
%  \inst{2}
%  National Graphene Institute,\\
%  University of Manchester}
%
\date{}

% }}*

% *{{ bibliography setup
\usepackage[backend=bibtex]{biblatex}
\bibliography{refs}
\renewcommand{\footnotesize}{\scriptsize}
\AtEveryCitekey{\clearfield{title}
                \clearfield{pagetotal}
                \clearfield{pages}}

% }}*

% *{{ graphics on front page
\titlegraphic{\vspace*{-1.5cm}\includegraphics[width=3cm]{figs/lan_logo.png}
\hfill\includegraphics[width=3cm]{figs/nownano_logo.png}}

% }}*

% *{{ titlepage + toc

% Section and subsections will appear in the presentation overview
% and table of contents.

\begin{document}

\begin{frame}[plain]
  \titlepage\end{frame}

%\begin{frame}{Outline}
%  \tableofcontents
%\end{frame}

% contents at ssubsection starts - comment to get rid
\AtBeginSection[]
{\begin{frame}<beamer>{Outline}
  \tableofcontents[currentsection]
  \end{frame}}

% }}*

% Setup: what's the question? EM in 2D (logarithms!)
% Monolayers: Keldysh, results, big things, word on metallics
% Bilayers (and beyond...?): potentials, approximations, exactitudes,
% experiments, decay rates?

% *{{ What is my question?

% define charge complexes, motivate their study, talk about why 2D materials
% are interesting.

\begin{frame}{What are charge complexes, why study them?}

\begin{itemize}
  \item Charge-carrier complexes are \blue{bound states} of two or more charges in a host
  material.
  \item Common example is an exciton.

  \begin{figure}[H]
    \floatbox[{\capbeside\thisfloatsetup{capbesideposition={left, center},
    capbesidewidth=0.5\textwidth}}]{figure}[\FBwidth]
    {\caption{Schematic view of an exciton - a bound state comprised of one
    electron and one hole.}\label{fig:exc_schematic}}
    {\includegraphics[width=0.5\textwidth]{figs/exciton_setup.pdf}}
  \end{figure}

  \item At low enough temperatures, charge complexes {\bf dominate} the optical
  response of semiconductors.

\end{itemize}

\end{frame}

\begin{frame}{Effective mass approximation}

  \begin{itemize}

    \item In conventional semiconductors,\footnote{Take III-Vs as an explicit
    example.} the binding of excitons may be described in the \blue{effective mass
    approximation}.

    \item The electronic band structures of these materials serve
    to supply effective masses for electrons and holes.

    \item We then solve Schr\"{o}dinger equations analogous to those for the
    hydrogen atom:

    \begin{equation}
      \left[-\sum_i\dfrac{\hbar^2}{2m_i}\nabla^2_i + \sum_{i<j}V({\bf
      r}_{ij})\right]\Psi(\{{\bf r}_i\}) = E\Psi(\{{\bf r}_i\})
    \end{equation}

  \end{itemize}
\end{frame}

% }}* end my question

% *{{ 2D EM, our monolayer model, QMC

\begin{frame}{Electromagnetism in Flatland}
\begin{itemize}
  \item If we lived in flatland (bona fide 2D space), electromagnetism would
  be a \blue{very different beast}.
  \item Poisson's equation would look the same, however, it's (two-dimensional)
  solution would predict a ``logarithmic Coulomb potential'' between pairs of
  point charges.\footnote{\
  Explanation for technical people: look at the inverse Fourier transform of
  the Green's function $G({k})=1/k^2$ in 2D and 3D. By power counting you can
  see that this is the case. EM in 1D is ever stranger.}
\end{itemize}
\begin{equation}
  \nabla^2 \phi({\bf r}) = -\dfrac{\rho({\bf r})}{\epsilon_0}.
\end{equation}
\begin{equation}
  V_q({\bf r}) \propto -q\ln{\lvert {\bf r} \rvert}.
\end{equation}
\begin{itemize}
  \item But we don't live in flatland, so we \red{don't get this}
  \frownie\
  \item We \green{do} have to consider screening, however.
\end{itemize}
\end{frame}

% \begin{frame}{Reality: 3D space, and electrostatic screening}
%   \ldots or do we? In any realistic model, and in every experiment ever
%   done on 2D materials, one other factor is very important: \green{screening}.
%
% \begin{itemize}
%   \item Electrostatic screening is the suppression in electric field strength
%   due to the presence of mobile charges.
%   \item Screening is lessened for 2D materials that are taken from their parent
%   compound - interactions between charges are stronger.
% \end{itemize}
% \end{frame}
%
\begin{frame}{The Keldysh Interaction}
\begin{itemize}
  \item Keldysh\footfullcite{Keldysh1979} showed that the effective interaction
  between charges in a polarisable 2D material has a special form
  \begin{equation}
    V_{\text{K}}({\bf r}) = \dfrac{q_1q_2}{8\epsilon_r \epsilon_0 r_*}
    \left[ H_0\left(\frac{r}{r_*}\right) - Y_0\left(\frac{r}{r_*}\right)
    \right],
  \end{equation}
  with $r_*$ a material parameter (usually \red{inferred} from \textit{ab initio}
  DFT or $GW$ calculations of $\epsilon({\bf q},\omega)$).
  \begin{figure}[H]
    \floatbox[{\capbeside\thisfloatsetup{capbesideposition={right, center},
    capbesidewidth=0.5\textwidth}}]{figure}[\FBwidth]
    {\caption{The Keldysh interaction. Notice we have recovered a
    log-dependence at short-range\ldots}\label{fig:keldysh}}
    {\includegraphics[width=0.5\textwidth]{figs/full_int.eps}}
  \end{figure}

\end{itemize}
\end{frame}

\begin{frame}{Complexes in Monolayers}
\begin{itemize}
  \item It turns out that the Keldysh interaction permits numerous bound states
  to exist, for typical $r_*$ values ($\mathcal{O}$(50\ \AA), vacuum).
\end{itemize}
\begin{figure}[H]
  \centering
  \includegraphics[width=0.8\linewidth]{figs/various_complexes.pdf}
  \caption{A few of the charge carrier complexes which can exist in 2D
  semiconductors.}\label{fig:various_complexes}
\end{figure}
\end{frame}


\begin{frame}{Our Model}{\& some of its limitations}
\begin{itemize}
  \item Taking monolayers as a goal, for now, we seek to solve
  \begin{equation}
      \left[-\sum_i\dfrac{\hbar^2}{2m_i}\nabla^2_i + \sum_{i<j}V_{\text{K}}({\bf
      r}_{ij})\right]\Psi(\{{\bf r}_i\}) = E\Psi(\{{\bf r}_i\}),
  \end{equation}
  for various complexes. We then compare total energies $E$ in order to
  evaluate binding energies (experimentally relevant).

  \item We ignore exchange effects, but can add
  them back in as a \blue{perturbative correction} (requires pair distribution
  functions, and unknown parameters\ldots).

  \item We are assuming that both the E.M.A. holds, and that the Keldysh
  interaction is a good approximation for real materials.
\end{itemize}
\end{frame}


\begin{frame}{Solving the few-body problem:\\ Quantum Monte Carlo}
\begin{itemize}
  \item We solve our few-body effective mass Schr\"{o}dinger equations by using
  the variational and diffusion quantum Monte Carlo methods (VMC, DMC).
  \item In VMC, estimates of high-dimensional integrals are formed from
  (cleverly weighted) random sampling. \red{The results are as good as the trial
  wave function}.
\end{itemize}
\begin{align}
  E\left[\ket{\Psi}\right] &= \dfrac{\bra{\Psi} \mathcal{\hat H}
  \ket{\Psi}}{\bra{\Psi}\Psi \rangle} = \displaystyle\int \dd{\bf R} \Pi({\bf
  R}) E_{\text{L}}({\bf R}) \nonumber \\
  \Pi({\bf R}) &= \dfrac{\lvert\Psi({\bf R}) \rvert^2}{\int\dd{\bf R}\lvert
  \Psi({\bf R})\rvert^2},\ E_{\text{L}} = \dfrac{\mathcal{\hat H} \Psi({\bf R})}{\Psi({\bf R})}
\end{align}
\end{frame}


\begin{frame}{Quantum Monte Carlo II - VMC}
\begin{itemize}
  \item VMC is usually only ever done as a prelude to DMC\@.
  \item In these models, we start with \blue{educated guesses} at the trial wave
  function which are of \blue{Jastrow} form
  \begin{equation}
    \Psi_{\text{T}}({\bf R}) = \exp{\left[ \mathcal{J}_{\{\alpha\}}({\bf R})\right]}
  \end{equation}
  where $\{\alpha\}$ are a set of optimisable parameters.
  \item In practice, we vary the $\{\alpha\}$ such that some property of the
  wave function is \blue{optimal} (minimise E, variance of E, or MAD of E).
\end{itemize}
\end{frame}

\begin{frame}{Quantum Monte Carlo III - DMC}
\begin{itemize}
  \item In DMC, a trial function is propagated in imaginary
  time,\footnote{Details not important - but we do this by interpreting the
  problem as a \textit{statistical} one.} such that
  any excited state components it may contain are removed. The results are
  \green{independent of the trial function (in our cases), and offer improvement over
  VMC results.}
  \item We exploit (with $t=i\tau$)
  \begin{equation}
    \lim_{\tau \rightarrow \infty} \exp{\left[ -\tau \mathcal{\hat H} \right]}
    \ket{\Psi_{\text{T}}} \sim \ket{\Psi_{\text{GS}}}
  \end{equation}

  \item DMC is \green{exact for nodeless wave functions}.

\end{itemize}
\end{frame}

\begin{frame}{Fairness: How else might one do this?}
Other approaches for solving these kinds of problems exist. I won't
talk about these, but for completeness I will mention them.
\begin{itemize}
  \item $GW$--$BSE$: Solve the Bethe-Salpeter equation for the spectral
  function of a material including two-body effects.
  \item First-principles QMC:\@ Solve the many-electron Schr\"{o}dinger
  equation for the interacting system. Few people have ever actually done this
  (I'm one!).
\end{itemize}
There's a good book covering both of these by Martin, Reining, and
Ceperley.\footfullcite{Martin2016}
\end{frame}


\begin{frame}{Results: Monolayers}

\begin{table}[H]
  \centering
\begin{tabular}{lcccccc}
\hline \hline
& \multicolumn{2}{c}{$E_{\rm X^-}^{\rm b}$
  (meV)} & \multicolumn{2}{c}{$E_{\rm
    X^+}^{\rm b}$ (meV)} \\
\raisebox{1.5ex}[0pt]{TMDC} & DMC&
Exp. &
DMC & Exp. \\ \hline

MoS$_2$ & $35.0$ &
$40$, $18.0(15)$, $43$
& $34.9$ & \\

MoSe$_2$ & $34.5$ &
$30$ & $34.4$ &
$30$ \\

WS$_2$ & $33.5$ &
$34$,$36$,
$10$--$15$, $30$,
$45$ & $33.6$ & & \\

WSe$_2$ & $29.6$ &
$30$ & $29.6$ &
$30$, $24$ \\

\hline \hline
\end{tabular}
\caption{Some trion binding energy results from our paper. Refs.\ and comparative results may be found
therein.\footfullcite{biexpaper}}\label{tab:trion_results}
\end{table}

\begin{itemize}
  \item We also enjoy good agreement with other theoretical studies, based on
  DMC, PIMC, SVM, heavy-hole approx., and variational calculations.
\end{itemize}

\end{frame}

% \begin{frame}{Results: Monolayers II}
%
% \begin{itemize}
%   \item Another aspect of this paper was the study of higher
%   complexes. Surprisingly, we find that such systems are bound\ldots
%
% \begin{table}[H] \caption{Binding energies of BIG complexes in different
% TMDCs. Note: these aren't necessarily PL peak
% positions.\label{table:big_complexes}}
%
% \begin{center}
% \begin{tabular}{lccc} \hline
% \hline
% % peak position == de-excitonisation energies.
%
%  & \multicolumn{3}{c}{Binding energy (meV)} \\
%
% \raisebox{1.5ex}[0pt]{TMDC} & XX$^-\ (\rm eeehh)$ & D$^-$X ($\rm Deeeh$) & D$^0$XX
% ($\rm Deeehh$) \\
%
% \hline
%
% MoS$_2$  & $58.6(6)$ & $84.4(4)$ & $61.6(6)$ \\
%
% MoSe$_2$ & $57.0(4)$ & $57.9(2)$ & $56.9(9)$ \\
%
% WS$_2$   & $57.4(3)$ & $59.2(4)$ & $58.2(6)$ \\
%
% WSe$_2$  & $52.5(7)$ & $51.3(4)$ & $51(1)$ \\
%
% \hline \hline
% \end{tabular}
% \end{center}
% \end{table}
%
% \item Biex.\ problem: Big things end up being called ``biexcitons''?
%
% \end{itemize}
%
% \end{frame}

% }}*

% *{{ Modelling heterostructures: multilayer interactions
\begin{frame}{Heterostructures: multilayer interactions}
\begin{itemize}
  \begin{figure}[H]
    \floatbox[{\capbeside\thisfloatsetup{capbesideposition={left, center},
    capbesidewidth=0.5\textwidth}}]{figure}[\FBwidth]
    {\caption{A tri-layer heterostructure of 2D semiconductors. There are $N$
    layer potentials, and $N(N-1)/2$ inter-layer potentials.}\label{fig:multilayers}}
    {\includegraphics[width=0.5\textwidth]{figs/multilayers.pdf}}
  \end{figure}

  \item Fourier components can be
  determined,\footnote{Analytically, upto $N=4$. Possibly higher $N$ in
  special cases, or, as always, with symmetry.} but
  real-space potentials obtained numerically (Hankel transform).
\end{itemize}
\end{frame}

\begin{frame}{Bilayers of 2D semiconductors}
How do charges interact in bilayers? More special functions?
\begin{figure}[H]
  \floatbox[{\capbeside\thisfloatsetup{capbesideposition={left, center},
  capbesidewidth=0.5\textwidth}}]{figure}[\FBwidth]
  {\caption{Exact inter-layer and intra-layer interaction potentials versus
  distance. Parameters taken from an experimentally relevant MoSe$_2$/WSe$_2$
  geometry.}\label{fig:interactions}}
  {\includegraphics[width=0.45\textwidth]{figs/both_wse2mose2.eps}}
\end{figure}

\begin{itemize}
  \item The screened interactions start to get cumbersome:
\end{itemize}
\begin{equation}
  v_{\text{intra}}({\bf q}) =
  \frac{(1+r_{*j}q)\exp(qD)
  -r_{*j}q\exp(-qD)}{2\bar{\epsilon}
  q\left[(1+r_{*j}q)(1+r_{*i}q)\exp(qD)-
  r_{*i}r_{*j}q^2\exp(-qD)\right]} \label{eq:intralayer_pot}
\end{equation}
\end{frame}

% \begin{frame}{A Quick Tangent}
%
% \begin{itemize}
%   \item The interactions are qualitatively similar to the (displaced) Coulomb
%   interactions. We expect similar physics to that in CQWs of III-Vs.
%
%
% \begin{figure}[H]
%   \floatbox[{\capbeside\thisfloatsetup{capbesideposition={left, center},
%   capbesidewidth=0.5\textwidth}}]{figure}[\FBwidth]
%   {\caption{Schematic of an indirect trion in a coupled quantum well.}\label{fig:trion_cqw}}
%   {\includegraphics[width=0.35\textwidth]{figs/trion_cqw.pdf}}
% \end{figure}
%
%   \item These were originally studied as the ``ideal testbed'' for finding BECs
%   of excitons. We've studied them before.~\footfullcite{trion_paper}
%
% \end{itemize}
%
% \end{frame}

\begin{frame}{Bilayers of 2D semiconductors: Results}
\begin{itemize}
  \item We have studied various complexes in various dielectric environments
  (act to renormalize $r_\star$).
  \item Are able to explain positions and widths (see later) of three candidate
  peaks relevant to a recent experiment.
\end{itemize}


\begin{table}[H]
  \centering
  \begin{tabular}{c|c}
  Complex & Binding Energy (meV) \\ \hline
  $\tilde{\rm X}$ & 103.95863(8) \\
  ${\rm D}^0_{\rm e}$ & 163.247871(3)\\
  $\tilde{\rm X}{\rm e} \rightarrow \tilde{\rm X}$ + {\rm e} & 4.226(4) \\
  ${\rm D}^0_{\rm e} \tilde{\rm X}\rightarrow {\rm D}^{0}_{\rm e} + \tilde{\rm X}$ & 11.41(3)
  \end{tabular}
  \caption{Binding energies of various complexes of interest in an
  hBN/MoSe$_2$/WSe$_2$/hBN system.}\label{tab:hbn_hbn_4.0_results}
\end{table}
\end{frame}


\begin{frame}{What else can QMC do for us?}

\begin{itemize} \item Whilst we're still on the topic of model systems, {\bf is
there anything useful QMC can do beyond calculate binding energies?}

  \item In collaboration with Mark Danovich, David Ruiz-Tijerina, and Volodya
  Fal'ko, we've\footnote{Myself, Neil Drummond and Marcin Szyniszewski.} used
  QMC + pen-and-paper to form estimates of {\it lifetimes\/} in perturbation
  theory.
  \end{itemize}

\begin{minipage}[t]{0.49\textwidth}
  \begin{figure}[H]
    \centering
    \includegraphics[width=\linewidth]{figs/lifetime.pdf}\label{fig:lifetime}
  \end{figure}
\end{minipage}%
\hfill
\begin{minipage}[t]{0.49\textwidth}
\vspace*{-.5cm}
\begin{equation}
    \frac{\left| \int \psi (\mathbf{r}, \mathbf{r}, \mathbf{0})
    \dd^2 \mathbf{r} \right|^2}{\int | \psi |^2 \dd^6 \mathbf{R}} \sim 10^{-7}
    \text{a.u.}^{-4} \nonumber
\end{equation}
\begin{itemize}
  \item $\psi$ the donor-bound trion (${\rm D}^0\tilde{\rm X}$).
\end{itemize}
\end{minipage}%
\end{frame}

% REPLACE W/ PHOSPHORENE CALCS?
\begin{frame}{First-principles QMC\@: Phosphorene}

\begin{block}{First-Principles Excited State QMC}
For a $n\times m$ supercell of phosphorene unit
cells, ground and (singlet) excited states could be modelled by:
\begin{align}
  \Psi({\bf R}) &= \exp{\left[ \mathcal{J}({\bf R}) \right]} \times
  \mathcal{D}_{\uparrow}({\bf R}) \times \mathcal{D}_{\downarrow}({\bf R}),
  \nonumber \\
  \Psi^{\star}({\bf R}) &= \exp{\left[ \mathcal{J}({\bf R}) \right]} \times
  \mathcal{D}^{\star}_{\uparrow}({\bf R}) \times \mathcal{D}_{\downarrow}({\bf
  R}).
\end{align}

$\mathcal{D}^{\star}$ is a Slater determinant of single-particle orbitals which
has had one orbital replaced with an excited orbital.\footnote{The important
point in phosphorene is the $\Gamma$ point\ldots}
\end{block}
\end{frame}

\begin{frame}{First-principles QMC\@: cont.}
The \textit{excitonic} and \textit{quasiparticle} energy gaps are then:
\begin{align}
  \Delta_{\text{X}} &= E \left[ \Psi^{\star}_N \right] - E \left[ \Psi_N \right]
  \nonumber \\
  \Delta_{\text{Q}} &= E \left[ \Psi_{N+1} \right] + E\left[ \Psi_{N-1} \right]
  - 2 E\left[ \Psi_N \right]
\end{align}
\vspace*{-.5cm}
\begin{block}{But it's not\footnote{Or, \textit{wasn't}.} clear QMC can
do this ``cleanly''\ldots}
\begin{minipage}[t]{0.45\textwidth}
\begin{itemize}
  \item Periodic BCs
  \begin{itemize}
    \item Need to remove \red{finite-size} effects.
    \item \& balance resources\ldots
  \end{itemize}
\end{itemize}
\end{minipage}%
\hfill
\begin{minipage}[t]{0.45\textwidth}
\begin{itemize}
  \item Variational Principles
  \begin{itemize}
    \item Excitonic energies \red{don't} have them\ldots
    \item Quasiparticle energies \green{do}.
  \end{itemize}
\end{itemize}
\end{minipage}%
\end{block}

\end{frame}

\begin{frame}{First-principles QMC\@: cont.}
We find in phosphorene\footnote{But have also studied hBN, Bulk hBN, Si,
SiO$_2$, c-BN.} that we \green{can} do this sensibly.

\begin{figure}[H]
  \centering
  \includegraphics[width=.8\linewidth]{figs/phos.pdf}
  \caption{DMC quasiparticle and excitonic gaps for phosphorene. Their
  difference is the exciton binding energy, $E^{\rm B}_{\rm X}$.}
\label{fig:phos}
\end{figure}

\end{frame}

% }}*

% *{{ Future Research
\begin{frame}{Future (relevant) Research Avenues}

\begin{itemize}

  \item Electron-hole droplets in TMDCs? There's been recent interest in
  studying ``electron-hole droplets''.\footfullcite{Keldysh1986,Almand2014}
  These might show up in TMDC bilayers. I'd like to find out.\footnote{This
  requires ``Ewaldising'' the bilayer Keldysh interaction - i.e.\ making
  lattice sums of its $1/r$ component absolutely convergent. I've recently done
  this for monolayers.\ldots}

  \item Electron gases in metallic TMDCs? Lots of literature on 2DEGs in
  context of III-Vs, but presumably metallic 2D materials will one day be
  of interest. Requires periodic Keldysh interaction.

\end{itemize}
\end{frame}
% }}*

% *{{ thanks!
\begin{frame}{Credits}
C $\rightarrow$ charge-complexes (various) \\
X $\rightarrow$ first-principles excitations in QMC\@.
\begin{minipage}[t]{0.45\textwidth}
\begin{block}{Lancaster}
\begin{itemize}
  \item N. Drummond (C, X)
  \item M. Szyniszewski (C, X)
  \item O. Witham (C)
  \item E. Mostaani (C)
\end{itemize}
\end{block}

\begin{block}{Japan - \textit{JAIST}}
\begin{itemize}
  \item R. Maezono (X)
  \item G. Prayogo (X)
\end{itemize}
\end{block}

\end{minipage}%
\hfill
\begin{minipage}[t]{0.45\textwidth}
\begin{block}{UoM - \textit{NGI} }
\begin{itemize}
  \item M. Danovich (C)
  \item D. Ruiz-Tijerina (C)
  \item V. Fal'ko (C, X)
  \item V. Z\'{o}lyomi (X)
\end{itemize}
\end{block}

\begin{block}{Money}
Thanks to EPSRC and the Graphene NOWNANO CDT (who fund me).
\end{block}


\end{minipage}%
\end{frame}
% }}*

\end{document}

